\section{Main Theorem}
\label{sec:main-theorem}

\begin{theorem}[Conditional polynomial probabilistic dimension]
\label{thm:main}
Under Assumptions~\ref{assump:fixed-source-witnesses},
\ref{assump:universal-expected-success},
\ref{assump:constant-depth}, and~\ref{assump:robust-tube}, the unconditional
initialization-gate law $\mathcal P_{\rm gate}$, selected before every
$\mathcal D\in\Delta(\mathcal X)$ and $h\in\mathcal H$, satisfies
\[
\mathbb E_{\varphi\sim\mathcal P_{\rm gate}}\!\left[
\inf_{w\in\mathbb R^{d_{\rm path}}}
R_{\mathcal D,h}(w,\varphi)
\right]\leq\varepsilon+\delta_0.
\tag{18}\label{eq:p1-i3-eq018}
\]
Consequently,
\[
\operatorname{dc}_{\varepsilon+\delta_0}(\mathcal H)
\leq d_{\rm path}
\leq S^L
\leq S^{L_0},
\tag{19}\label{eq:p1-i3-eq019}
\]
and the same law is admissible at the larger threshold $2\varepsilon$, so
\[
\operatorname{dc}_{2\varepsilon}(\mathcal H)\leq S^{L_0}.
\tag{20}\label{eq:p1-i3-eq020}
\]
These bounds are nonasymptotic and have no hidden multiplicative constant.
The exposed quantitative variables are
$S,L,L_0,\varepsilon,\delta_0,T,\eta,r$ through the assumptions and the exact
conclusions above; $L_0$ is universally fixed.  The learner premise is an
expectation over initialization and the fixed finite sample horizon, the tube
premise is an initialization probability, and \eqref{eq:p1-i3-eq018} is an expectation over the
unconditional law $\mathcal P_{\rm gate}$.  The metric is the exact
tie-resolved classification risk, and no coefficient norm or positive score
margin is required.
\end{theorem}

For $L=1$, the feature map is the gate-free map $x\mapsto x$ and
$d_{\rm path}=S$.  If $\delta_0=0$, the complement penalty disappears and
\eqref{eq:p1-i3-eq018} has right-hand side $\varepsilon$.  If
$\varepsilon=\delta_0=0$, both public error thresholds remain exactly zero.
These reductions are part of the theorem rather than asymptotic limiting
claims.
